\documentclass[12pt]{article}
\usepackage[utf8]{inputenc}
\usepackage{graphicx}
\usepackage{hyperref}
\usepackage{amsmath}
\usepackage{geometry}
\usepackage{setspace}

% 1-inch margins and 1.5 spacing for readability and length
\geometry{a4paper, margin=1in}
\onehalfspacing

\title{\textbf{Deep Learning for 3D Bone Analysis:}\\ Building a Proximal Humerus Atlas and Fracture Map}
\author{Project Report}
\date{\today}

\begin{document}

\maketitle
\tableofcontents
\newpage

\begin{abstract}
This report details the development of a statistical shape model (atlas) and a fracture density map for the proximal humerus bone. Using a dataset of 30 fractured patient scans, we adapted the \textbf{CINeMA} (Conditional Implicit Neural Multi-Modal Atlas) framework—originally designed for brain MRI—to work with 3D bone fragment surfaces. The project involves a complete end-to-end pipeline: converting raw STL meshes into volumetric Signed Distance Fields (SDFs), training a SIREN-based neural network to learn a canonical bone shape, and aggregating fracture patterns into a population-wide density map. The final output is an interactive 3D tool that visualizes the "mean" humerus and identifies high-risk fracture zones, aiding in the design of orthopedic implants.
\end{abstract}

\section{Introduction}
Fractures of the proximal humerus (the top part of the arm bone) are common and complex. Designing surgical plates and nails requires a deep understanding of the "average" bone shape and typically where it breaks. 

Traditional methods often rely on manual measurements or simple averaging of healthy bones. Our approach uses \textbf{Implicit Neural Representations (INRs)}, a modern deep learning technique that represents 3D shapes as continuous mathematical functions rather than discrete pixels or meshes. This allows for infinite resolution and enables us to learn a smooth, "canonical" atlas even from sparse, fragmented input data.

\section{Dataset Description}

\subsection{The Raw Data}
The dataset consists of 30 distinct subjects, divided into two categories based on the \textit{quality and alignment} of the fragments:

\begin{itemize}
    \item \textbf{Ground Truth (GT) Reconstructions}: (e.g., \texttt{GT 02}, \texttt{GT 03})
    \begin{itemize}
        \item These folders contain bone fragments, just like the patient folders.
        \item \textbf{Crucial Difference}: In GT data, the fragments have been \textbf{perfectly aligned} by a clinical expert. They sit in 3D space exactly where they belong to form a complete humerus.
        \item Think of this as a \textbf{solved jigsaw puzzle}. We use these to train the neural network because they represent the "correct" shape.
    \end{itemize}
    
    \item \textbf{Patient Scans}: (e.g., \texttt{patient\_00})
    \begin{itemize}
        \item These folders also contain bone fragments.
        \item \textbf{Crucial Difference}: These are raw clinical scans where fragments might be displaced, rotated, or "exploded" due to the injury.
        \item Think of this as a \textbf{jumbled jigsaw puzzle}. We use the trained Atlas to help "solve" these puzzles by matching the fragments to the canonical shape.
    \end{itemize}
\end{itemize}

Every subject is provided as a collection of \textbf{STL (Stereolithography)} files. An STL file defines the surface of a 3D object using triangles. Unlike medical CT scans (which are 3D grids of pixels), these files only describe the \textit{shell} of the bone.

\subsection{The Challenge}
Neural networks cannot easily "eat" raw triangle meshes. They prefer structured grids (volumes). Furthermore, real-world data is messy:
\begin{enumerate}
    \item \textbf{Side Inconsistency}: Some bones are from the Left arm, some from the Right.
    \item \textbf{Positioning}: Every scan is in a different position and rotation in 3D space.
    \item \textbf{Fragmentation}: The "Patient" files are broken into pieces.
\end{enumerate}

\section{Methodology: Data Pipeline}

To solve the challenges above, we built a robust preprocessing pipeline (\texttt{mesh\_to\_volume.py} and \texttt{prepare\_data.py}).

\subsection{1. Virtual Surgery (Fragment Assembly)}
For patient scans, we first mathematically combine all individual fragment STLs into a single mesh object. This simulates "gluing" the bone back together for analysis.

\subsection{2. Bilateral Mirroring}
A Left humerus is essentially a mirror image of a Right humerus. To double our effective training data and prevent the model from learning two different shapes (Left vs. Right), we detect the side of each bone. If a bone is from the Left side, we mathematically flip it along the X-axis:
\begin{equation}
    x_{new} = -x_{old}
\end{equation}
This ensures the network only sees "Right-sided" humeri, making learning much faster and more accurate.

\subsection{3. Voxelization and SDF Calculation}
This is the most critical step. We convert the hollow surface mesh into a solid 3D volume using a \textbf{Signed Distance Field (SDF)}.
\begin{itemize}
    \item We create a 3D grid (128x128x128 voxels).
    \item For every point in this grid, we calculate the distance to the nearest point on the bone surface.
    \item \textbf{Inside} the bone, the distance is negative.
    \item \textbf{Outside}, the distance is positive.
    \item \textbf{On the surface}, the distance is zero.
\end{itemize}
This transforms the messy triangle soup into a smooth, continuous field that tells the network exactly where the bone is.

\subsection{4. Data Splitting}
We split the 30 subjects into two sets:
\begin{itemize}
    \item \textbf{Training Set (80\% = 24 subjects)}: Used to teach the network the shape of the humerus. The model sees these bones thousands of times.
    \item \textbf{Validation Set (20\% = 6 subjects)}: Kept hidden during training. We use these to test if the model can generalize to new bones it hasn't seen before.
\end{itemize}
This split is defined in \texttt{subjects.yaml} to ensure reproducibility.

\section{Methodology: The Neural Atlas (The "Brain")}

We adapted the \textbf{CINeMA} architecture. While the original CINeMA was built for brain MRI (handling multiple tissue types), we simplified it to focus purely on bone geometry.

\subsection{Network Architecture: SIREN}
We use a specialized neural network called a \textbf{SIREN} (Sinusoidal Representation Network). 
\begin{itemize}
    \item \textbf{Standard Networks}: Use ReLU activation functions (rectified linear units). These are good for classifying images (cat vs. dog) but struggle to represent fine 3D details and smooth curves.
    \item \textbf{SIREN}: Uses \textbf{sine waves} as activation functions.
    \begin{equation}
        \phi(x) = \sin(\omega_0 \cdot x)
    \end{equation}
    Sine waves are periodic and differentiable, allowing the network to capture complex, high-frequency details of the bone surface (like the sharp ridges of the tuberosities) with high precision.
\end{itemize}

\subsection{The Auto-Decoder Framework}
Typical neural networks take an input (image) and produce an output (label). Our "Auto-Decoder" works differently. It doesn't take an image as input. Instead, it learns a "code" (latent vector) for each patient.

\begin{itemize}
    \item \textbf{Latent Code ($z_i$)}: A descriptive vector that summarizes the unique shape of Patient $i$.
    \item \textbf{Decoder Network ($f_\theta$)}: A shared function that turns any code and a 3D coordinate $(x,y,z)$ into an SDF value.
    \item \textbf{Transformation ($\tau_i$)}: The model also learns how to rotate and translate the "mean" atlas to match Patient $i$'s specific position.
\end{itemize}

Training involves optimizing \textbf{everything at once}: the network weights ($\theta$), every patient's shape code ($z_i$), and every patient's position ($\tau_i$).

\section{Methodology: Fracture Density Mapping}

After training the atlas, we have a "Canonical Humerus"—the perfect average shape. But we also want to know where it breaks.

\subsection{Mapping Fractures to Atlas Space}
Since we learned a transformation for every patient (rotation and translation), we can invert this transformation.

\begin{enumerate}
    \item We take the original fracture fragments of Patient $i$.
    \item We identify the "cracks" (the boundaries between fragments).
    \item We apply the \textbf{inverse transformation} to move these cracks from "Patient Space" to "Atlas Space".
\end{enumerate}

\subsection{Building the Heatmap}
We do this for all 30 patients. Now, in the Atlas Space, we have 30 sets of ghost fractures floating around the mean bone.
\begin{itemize}
    \item We stack them all up in a 3D grid.
    \item We count how many patients define a fracture at each voxel.
    \item This creates a **3D Probability Cloud**.
    \item \textbf{Red Zones}: Areas where almost everyone broke their bone (High Probability).
    \item \textbf{Clear Zones}: Areas that rarely break (Low Probability).
\end{itemize}

\section{Experimental Results}

\subsection{Training Performance}
We trained the model for 50 epochs on an NVIDIA L4 GPU.
\begin{itemize}
    \item \textbf{Loss Convergence}: The error (difference between predicted and actual shape) dropped consistently, stabilizing around 0.01.
    \item \textbf{Time}: Training took approximately 1 hour.
\end{itemize}

\subsection{The Final Atlas}
The resulting "Mean Humerus" is smooth and anatomically accurate. It effectively averaged out the noise and partial fragments from the inputs to hallucinate a complete, healthy bone.

\subsection{Fracture Hotspots}
The generated density map (\texttt{fracture\_density.npy}) revealed two major critical zones:
\begin{enumerate}
    \item \textbf{Surgical Neck}: A ring around the bone just below the head. This is the most common site for "2-part" fractures.
    \item \textbf{Greater Tuberosity}: The side of the shoulder, often sheared off in "3-part" or "4-part" fractures.
\end{enumerate}
These findings perfectly align with orthopedic classification systems (Neer Classification), validating our automated approach.

\section{Conclusion}
This project successfully demonstrated that Deep Learning—specifically Implicit Neural Representations—can be used to build high-quality statistical shape models from fractured bone data. We provided a complete pipeline from raw STL files to a clinical visualization tool. The resulting Atlas and Fracture Map provide valuable quantitative data for improving surgical hardware, ensuring plates and screws are placed exactly where bones are most likely to fail.

\end{document}
